
\section{\BookRecordIndex{T0177}加权范数}\label{sec:10-3}

为了应用\BookRecordIndex{T0380}隐函数定理，必须指明 $\bar\partial_J(u^R)$ 在哪一种范数下很小。选择范数时的
指导原则，是曲线 $u^0$ 与 $u^\infty$ 在这个\BookRecordIndex{T0254}粘合论证中地位相同。然而，映射
$u^\infty$ 以重标度后的形式出现，并集中在半径 $1/R$ 的球内。因此，为赋予 $u^0$ 与
$u^\infty$ 相同的权重，我们考察 $2$-球面上一族依赖于 $R$ 的度量，使半径 $1/R$ 球的
体积近似等于 $S^2$ 关于标准 Fubini--Study 度量的体积。重标度后的度量形如
\begin{equation}\label{eq:10-3-1}
  g^R:=(\theta^R)^{-2}(ds^2+dt^2),
\end{equation}
其中
\[
  \theta^R(z):=
  \begin{cases}
    R^{-2}+R^2|z|^2,& |z|\le1/R,\\
    1+|z|^2,& |z|\ge1/R.
  \end{cases}
\]
相应的体积形式是
$\mathrm{dvol}^R:=(\theta^R)^{-2}\,ds\wedge dt$，而 $S^2$ 关于这个度量的面积为
\begin{equation}\label{eq:10-3-2}
  \operatorname{Vol}^R(S^2)
  =\int_{\mathbb C}(\theta^R)^{-2}\,ds\wedge dt
  \le2\pi.
\end{equation}

有两种方式可以理解这个度量。首先注意，对合
$z\mapsto1/(R^2z)$ 是关于 $g^R$ 的等距映射，并交换两个圆盘
$\{|z|\le1/R\}$ 与 $\{|z|\ge1/R\}$。因此，许多计算只须在外圆盘
$\{|z|\ge1/R\}$ 上进行；那里诱导的度量就是通常的 Fubini--Study 度量。考察重标度
$z\mapsto R^2z$ 也很有用。它给出从配备度量 $g^R$ 的圆盘 $\{|z|\le1/R\}$ 到配备
Fubini--Study 度量 $g^1$ 的圆盘 $\{|z|\le R\}$ 的等距映射。因此，这个度量对沿
$u^R$ 的向量场之 $L^p$ 和 $W^{1,p}$ 范数有如下作用。若把支集位于
$B_{\delta/(2R)}$ 中的向量场 $\xi(z)\in T_{u^R(z)}M$ 重标度为沿 $u^\infty$ 的向量场
\[
  \xi^\infty(z):=\xi(z/R^2)\in T_{u^\infty(z)}M,
\]
那么重标度后向量场 $\xi^\infty$ 的标准范数与原向量场 $\xi$ 的\BookRecordIndex{T0177}加权范数一致。

更明确地说，回忆\BookRecordIndex{T0088}殆复结构 $J^R=\{J_z^R\}_{z\in S^2}$ 由
式~\eqref{eq:10-1-4} 定义，并定义\BookRecordIndex{T0177}加权范数
\[
  \|\xi\|_{0,p,R}
  :=\left(
  \int_{\mathbb C}\theta^R(z)^{-2}|\xi(z)|_{J_z^R}^p\,ds\,dt
  \right)^{1/p},
\]
\[
  \|\xi\|_{1,p,R}
  :=\left(
  \int_{\mathbb C}
  \left(
    \theta^R(z)^{-2}|\xi(z)|_{J_z^R}^p
    +\theta^R(z)^{p-2}|\nabla\xi(z)|_{J_z^R}^p
  \right)ds\,dt
  \right)^{1/p}.
\]
这里 $\nabla=\nabla^{R,z}$ 表示由 $\omega$ 和 $J_z^R$ 经~\eqref{eq:2-1-1} 确定的
$M$ 上度量之 \BookRecordIndex{T0034}Levi--Civita 联络。此外，我们在 $\mathbb C$ 上使用坐标 $z=s+it$，并写
\[
  |\nabla\xi(z)|_J^2
  :=|\nabla_s\xi(z)|_J^2+|\nabla_t\xi(z)|_J^2.
\]
对
$\eta=\eta_1\,ds+\eta_2\,dt\in\Omega^1(S^2,(u^R)^*TM)$，写
\[
  |\eta|^2:=|\eta_1|^2+|\eta_2|^2,
  \qquad
  |\nabla\eta|_J^2
  :=|\nabla_s\eta_1|_J^2+|\nabla_t\eta_1|_J^2
    +|\nabla_s\eta_2|_J^2+|\nabla_t\eta_2|_J^2,
\]
并定义
\[
  \|\eta\|_{0,p,R}
  :=\left(
  \int_{\mathbb C}\theta^R(z)^{p-2}|\eta(z)|_{J_z^R}^p\,ds\,dt
  \right)^{1/p},
\]
\[
  \|\eta\|_{0,\infty,R}
  :=\sup_{z\in\mathbb C}\theta^R(z)|\eta(z)|_{J_z^R},
\]
\[
  \|\eta\|_{1,p,R}
  :=\left(
  \int_{\mathbb C}
  \left(
    \theta^R(z)^{p-2}|\eta(z)|_{J_z^R}^p
    +\theta^R(z)^{2p-2}|\nabla\eta(z)|_{J_z^R}^p
  \right)ds\,dt
  \right)^{1/p}.
\]

当 $R=1$ 且 $J_z^R\equiv J$ 时，这些就是把 $\mathbb C$ 看作
$\mathbb{CP}^1$ 的坐标片后，关于 Fubini--Study 度量的通常 $L^p$ 与 $W^{1,p}$ 范数。
对一般的 $R$，应分两部分理解这些范数。在区域 $|z|\ge1/R$ 中，它们仍是通常的范数；
而在区域 $|z|\le1/R$ 中，它们与沿 $u^\infty(z)$、位于半径 $R$ 球内的重标度向量场
\[
  \xi^\infty(z):=\xi(z/R^2)
\]
或 $1$-形式
\[
  \eta^\infty
  :=R^{-2}\bigl(\eta_1(z/R^2)\,ds+\eta_2(z/R^2)\,dt\bigr)
\]
关于 $\mathbb C$ 上 Fubini--Study 度量的通常范数一致。因此，如下引理所示，我们得到
常数与 $R$ 无关的通常索伯列夫估计。

\begin{lemma}\label{lem:10-3-1}
设 $A^0,A^\infty\in H_2(M;\mathbb Z)$，并给定满足~\eqref{eq:10-1-1} 的
$J^0,J^\infty\in\mathcal J(S^2;M,\omega)$。固定常数 $c>0$ 和 $p\ge1$，其中
$p\ne2$。则存在常数 $c_0=c_0(p,c)>0$，使得对所有满足~\eqref{eq:10-2-3} 的
$\delta,R$ 以及所有 $u=(u^0,u^\infty)\in\mathcal M(c)$，下列结论成立；其中
$u^R=u^{\delta,R}:S^2\to M$ 由~\eqref{eq:10-2-1}、\eqref{eq:10-2-2} 给出。

\begin{enumerate}[label=\textnormal{(\roman*)}]
\item 若 $p>2$，则
\[
  \|\xi\|_{L^\infty}\le c_0\|\xi\|_{1,p,R},
  \qquad
  \|\eta\|_{0,\infty,R}\le c_0\|\eta\|_{1,p,R},
\]
其中
$\xi\in\Omega^0(S^2,(u^R)^*TM)$，
$\eta\in\Omega^1(S^2,(u^R)^*TM)$。

\item 若 $1\le p<2$ 且 $1\le q\le2p/(2-p)$，则
\[
  \|\xi\|_{0,q,R}\le c_0\|\xi\|_{1,p,R},
  \qquad
  \|\eta\|_{0,q,R}\le c_0\|\eta\|_{1,p,R},
\]
其中
$\xi\in\Omega^0(S^2,(u^R)^*TM)$，
$\eta\in\Omega^1(S^2,(u^R)^*TM)$。
\end{enumerate}
\end{lemma}

\begin{proof}
两个断言都分别在区域 $\{|z|\ge1/R\}$ 和 $\{|z|\le1/R\}$ 上证明。先考察区域
$\{|z|\ge1/R\}$，并利用其中的度量是 $S^2\setminus B_{1/R}$ 上的
Fubini--Study 度量（因而与 $R$ 无关）这一事实。对 $p>2$，函数的结论由附录 B 的结果
得出。每一点 $z_0\in S^2\setminus B_{1/R}$ 都包含在
$S^2\setminus B_{1/R}$ 内某个半径为 $\pi/4$ 的圆盘 $D$ 中。所有这样的圆盘彼此等距；
又因为 $d_{\mathrm{FS}}(z,0)=\tan^{-1}(|z|)$，它们可与
$(\mathbb C,g^1)$ 中的单位圆盘 $B_1(0)$ 等同。这个度量与欧氏度量等价，所以
结论由引理 B.1.16 和命题 B.1.9 立即得出。向量场 $\xi$ 与 $1$-形式 $\eta$ 的相应结论
可仿照注~\ref{rem:3-5-1} 证明。

在 $p<2$ 的情形，考察引理 B.1.17 的证明。该证明表明：对单位球外为零的函数，沿着从
单位球外出发、平行于实轴和虚轴的射线积分，可在区域 $\{|z|\ge1/R\}$ 上得到常数一致的
估计。现在用\BookRecordIndex{T0188}截断函数作分解：一部分位于 $|z|=1/R$ 附近且支集包含在 $B_1$ 中，另一部分
的支集包含在 $S^2\setminus B_{1/2}$ 中；由此便得到整个
$S^2\setminus B_{1/R}$ 上的一致估计。前面的论证也适用于支集包含在
$S^2\setminus B_{1/2}$ 中的函数，因为这个区域同样可以等同于 $\mathbb C$ 中以原点为
中心的一个球，其度量与欧氏度量等价。

这证明了区域 $\{|z|\ge1/R\}$ 上的估计。在区域 $\{|z|\le1/R\}$ 中，结论由映射
$u^R$ 和\BookRecordIndex{T0291}权函数 $\theta^R$ 定义公式中的对合对称性得出（见下面的习题~\ref{ex:10-5-2}）。
引理得证。
\end{proof}

下一个引理在我们应用\BookRecordIndex{T0380}隐函数定理（定理~\ref{thm:3-5-2}）时起关键作用。后一定理的陈述
方式使它适用于如下情形：黎曼曲面 $\Sigma=\mathbb{CP}^1$ 上的度量具有一致有界的
索伯列夫嵌入常数，而光滑映射 $u:\Sigma\to M$ 的一阶导数具有一致的 $L^p$ 上界。
引理~\ref{lem:10-3-1} 断言度量 $(\theta^R)^{-2}(ds^2+dt^2)$ 满足第一个假设；下面的
引理~\ref{lem:10-3-2} 则断言\BookRecordIndex{T0227}连通和 $u^R$ 满足第二个假设。其证明方法是以后更精细计算的
范例。

\begin{lemma}\label{lem:10-3-2}
设 $A^0,A^\infty$ 和 $J^0,J^\infty$ 如引理~\ref{lem:10-3-1}。固定两个常数
$c>0$ 与 $p>2$。则存在常数 $c_0=c_0(c,p)>0$，使得下列结论成立。若 $\delta,R$
满足~\eqref{eq:10-2-3}，则对每个 $(u^0,u^\infty)\in\mathcal M(c)$，由
\eqref{eq:10-2-1}、\eqref{eq:10-2-2} 定义的函数 $u^R$ 都满足估计
\[
  \|du^R\|_{0,\infty,R}\le c_0,
  \qquad
  \|du^R\|_{0,p,R}\le c_0,
  \qquad
  \|\bar\partial_{J^R}(u^R)\|_{0,p,R}
  \le c_0(\delta R)^{-2/p}.
\]
\end{lemma}

\begin{proof}
前两个估计的证明是容易的习题。它利用了式~\eqref{eq:10-2-2} 关于 $u^0,u^\infty$ 的
对称性（这将在下面的习题~\ref{ex:10-5-2} 中明确写出），以及如下事实：由于 $\zeta^0(0)=0$，
在区域 $|z|\le2/(\delta R)$ 中，乘积 $\delta R\zeta^0(z)$ 一致地由一个与
$\delta,R$ 及 $(u^0,u^\infty)$ 无关的常数控制。

为证明第三个估计，令 $r:=\delta R$，并回忆注~\ref{rem:10-2-2}：当 $|z|\ge1/R$ 时
$u^R(z)=u^{0,r}(z)$，而在 $r\to\infty$ 时，$u^{0,r}$ 在 $W^{1,p}$ 范数下收敛到
$u^0$。再次利用对称性可得
\[
  \|\bar\partial_{J^R}(u^R)\|_{0,p,R}\longrightarrow0.
\]
为得到更精确的收敛速率估计，注意当 $|z|\le2/r$ 且 $r=\delta R$ 充分大时（见
式~\eqref{eq:10-2-2}），
\[
  u^{0,r}(z)=\exp_x(\rho(rz)\zeta^0(z)),
  \qquad
  \exp_x(\zeta^0(z))=u^0(z).
\]
因此，项 $\bar\partial_{J^0}(u^{0,r})$ 的支集包含在\BookRecordIndex{T0170}环域
$B_{2/r}\setminus B_{1/r}$ 中，并且在该区域上可写成两项之和。一项含有\BookRecordIndex{T0188}截断函数
$\rho(rz)$ 的一阶导数（其增长阶为 $r$）乘以函数 $\zeta^0(z)$；当 $|z|\le2/r$ 时，
后者被常数乘 $1/r$ 控制，故该乘积一致有界。另一项含有 $\zeta^0(z)$ 的一阶导数，它们
也一致有界。由于这个估计中的积分区域面积至多为 $4\pi r^{-2}$，引理得证。
\end{proof}

\begin{exercise}\label{ex:10-3-3}
本习题说明命题~\ref{prop:3-1-1} 中定义的算子 $D_{u^R}$ 一致有界。设
$A^0,A^\infty\in H_2(M;\mathbb Z)$，并取
$(J^0,J^\infty)\in\mathcal J_{\mathrm{reg}}(S^2;A^{0,\infty})$。固定两个常数
$c>0$ 与 $p>2$。证明存在常数 $c_0$，使下列结论对所有满足~\eqref{eq:10-2-3} 的
$\delta,R$ 都成立：若 $(u^0,u^\infty)\in\mathcal M(c)$，且
$u^R:S^2\to M$ 由~\eqref{eq:10-2-2} 定义，则对每个
$\xi\in\Omega^0(S^2,(u^R)^*TM)$，都有
\[
  \|D_{u^R}\xi\|_{0,p,R}\le c_0\|\xi\|_{1,p,R}.
\]
\end{exercise}

\begin{remark}\label{rem:10-3-4}
固定一对 $(u^0,u^\infty)\in\mathcal M(A^{0,\infty};J^{0,\infty})$。对每对满足
\eqref{eq:10-2-3} 的实数 $\delta,R$，用~\eqref{eq:10-2-2} 定义
$u^R:S^2\to M$。于是每个算子 $D_{u^R}$ 都满足形如
\begin{equation}\label{eq:10-3-3}
  \|\xi\|_{k+1,p,R}
  \le c_{k,p,\delta,R}
  \bigl(\|D_{u^R}\xi\|_{k,p,R}+\|\xi\|_{k,p,R}\bigr)
\end{equation}
的估计，其中 $\xi\in\Omega^0(S^2,(u^R)^*TM)$。此时可以问：常数
$c_{k,p,\delta,R}$ 是否能选得与 $R$ 无关？值得注意的是，$k=0$ 时答案为是，而
$kp>2$ 时答案为否。

$k=0$ 时的结论很不平凡。先考察 $u^0,u^\infty$ 都是常值的情形，再通过
引理~\ref{lem:10-4-2} 对内边界上一点处消失的\BookRecordIndex{T0170}环域函数所作的一项改进，把一般情形归约到
这一情形，便可由下面的命题~\ref{prop:10-5-1} 得出结论。因为我们的\BookRecordIndex{T0254}粘合论证审慎地避开了
估计~\eqref{eq:10-3-3}，这里不展开细节。

为考察 $kp>2$ 的情形，考虑两条常值曲线 $u^0(z)\equiv u^\infty(z)\equiv x$。于是
\BookRecordIndex{T0395}预粘合曲线 $u^R$ 也为常值，$D_{u^R}$ 是从 $S^2$ 到 $T_xM$ 的函数之标准
\BookRecordIndex{T0210}柯西--黎曼算子。因此，问题在于\BookRecordIndex{T0291}权函数 $\theta^R$ 对 \BookRecordIndex{T0012}Calderón--Zygmund 不等式
有何影响。如下形式的截面给出一个富有启发性的例子：
\[
  \xi^R(z):=\beta^R(z)\left(z+\frac{1}{R^2z}\right)\xi_0,
\]
其中 $\xi_0\in T_xM$ 是非零向量，而乘以 $i=\sqrt{-1}$ 要通过点 $x\in M$ 处的\BookRecordIndex{T0088}殆复
结构 $J(x)$ 来理解。假设 $\beta^R:\mathbb C\to[0,1]$ 是光滑\BookRecordIndex{T0188}截断函数，在区域
$|z|\ge1$ 中与 $R$ 无关，并满足
\[
  \beta^R(z)=\beta^R\left(\frac{1}{R^2z}\right)
  =
  \begin{cases}
    1,& 1/R^2\le|z|\le1,\\
    0,& |z|\ge2\text{ 或 }|z|\le1/(2R^2).
  \end{cases}
\]
由对称性，只须考察区域 $|z|\ge1/R$ 中估计式的各项。函数
$z\mapsto z+1/(R^2z)$ 是全纯的，且其各阶导数在区域 $1\le|z|\le2$ 中一致有界，故
\[
  \sup_R\|\bar\partial\xi^R\|_{k,p,R}<\infty
\]
对每个 $k$ 和每个 $p$ 都成立。$\xi^R$ 的 $(k,p,R)$-范数由其在 $|z|=1/R$ 附近的
行为决定；在这里，项 $z\mapsto1/(R^2z)$ 在区域 $1/R\le|z|\le1$ 中占主导。简单计算
表明，只要 $(k-1)p<2$，$\xi^R$ 的 $(k,p,R)$-范数就一致有界。同一计算还表明，对每个
满足 $kp>2$ 的整数 $k$，存在常数 $c=c(k,p)>0$，使得
\[
  c^{-1}R^{k-2/p}
  \le\|\xi^R\|_{k+1,p,R}
  \le cR^{k-2/p}.
\]
这意味着，只要 $k\ge1$ 且 $kp>2$，式~\eqref{eq:10-3-3} 中的常数
$c_{k,p,\delta,R}$ 就不能选得与 $R$ 无关。
\end{remark}
